\begin{theorem}[Parent kernel within the full regular closure span]%
    \label{thm:peps_regular_closure_kernel_intersection}
    \lean{TNLean.PEPS.torusPlaquetteHolonomy_closure_corner_eq_one_iff_commute,
        TNLean.PEPS.IsGInjective.exists_torusClosureCommutationFilter,
        TNLean.PEPS.IsGInjective.inf_parentKernel_span_torusGClosureClass_eq_commutingSpan,
        TNLean.PEPS.IsGInjective.finrank_parentKernel_inf_closureSpan}
    \leanok
    \uses{lem:peps_regular_region_holonomy_filter,
        thm:peps_regular_closure_parent_constraint, thm:peps_pair_centralizer_count}
    Let both native torus periods be at least three, and let the
    tensor be regular $G$-injective. Let $\mathcal C$ be the span
    of all native pair-conjugacy-class closure vectors, including
    noncommuting labels, and let $\mathcal C_{\mathrm{comm}}$ be
    the commuting closure span. Let $H$ be a finite sum of positive
    original parent interactions on connected regions whose
    actual closed-cell realizations are simply connected. Assume
    that this family contains the plaquette at the intersection
    of the two closure seams. Then
    \[
        \ker H\cap\mathcal C=\mathcal C_{\mathrm{comm}}.
    \]
    Its dimension is the number of commuting pair-conjugacy
    classes, equivalently the centralizer sum.
    This equality applies to arbitrary coherent sums of closure
    vectors and requires no $G$-isometry. It proves the final
    closure constraint within the finite closure span in the
    regular case of \cite[Theorem~5.5]{Schuch2010PEPS}.
    The separate assertion that every parent-kernel vector lies
    in $\mathcal C$ remains unproved.
\end{theorem}

\begin{proof}\leanok
    At the seam-intersection plaquette, identity holonomy is
    equivalent to commutativity of the two closure labels. The
    identity-holonomy filter therefore fixes every commuting
    closure and removes every noncommuting closure. Its image
    of $\mathcal C$ is contained in
    $\mathcal C_{\mathrm{comm}}$. Positivity of the parent sum
    forces each of its local constraints separately, so every
    vector in $\ker H$ is fixed by the corner filter. Such a
    vector in $\mathcal C$ thus belongs to
    $\mathcal C_{\mathrm{comm}}$. The reverse inclusion is the
    local parent inclusion for commuting closures.
\end{proof}
